% Theorem 1' : direction-indexed spatial identification for an embedding-valued
% exposure. Strengthens Theorem 1 from a fixed scalar treatment to a data-driven
% linear summary of a high-dimensional spatially structured field. Verified
% symbolically (theorem_1prime_verify.py) and numerically with flexible-ML DML
% (theorem_1prime_numerical.py).

\subsection{Identification is a property of the direction, not of the embedding}
\label{sec:thm1prime}

The treatment in a text-augmented valuation is not handed to the analyst as a
scalar. It is a direction chosen inside a high-dimensional representation, and the
same embedding admits a continuum of such directions, each a candidate summary of
what the listing says. Whether the resulting effect is identified turns out to
depend on which direction is chosen, and Theorem~\ref{thm:idprime} makes that
dependence exact. The scalar identification result of
Section~\ref{sec:identification} is the one-dimensional shadow of it.

Let $W(\ell)\in\mathbb{R}^p$ be the residualized embedding at location $\ell$,
carrying a smooth spatial channel that a coordinate control $L$ recovers, a
non-smooth spatial channel that discrete identifiers such as neighborhood and
school names load onto and that $L$ misses, and a location-orthogonal semantic
channel. Write the projection of each channel onto a candidate unit direction $v$
as $C(v)$, $D(v)$, and $E(v)$, so that the treatment $T=v'W$ decomposes as
$C(v)$ times the smooth channel plus $D(v)$ times the missed channel plus $E(v)$
times the semantic channel. The confounder $U=h(\ell)$ places weight $b$ on the
missed channel, and the outcome follows the partially linear form
$Y=\theta T+\beta U+\varepsilon$.

\begin{theorem}[Direction-indexed identification]
\label{thm:idprime}
Under Assumptions~\ref{as:id-U}--\ref{as:id-exog}, the Robinson estimand of the
effect of $T=v'W$ that conditions on $L$ equals
\begin{equation}
\theta(v)\;=\;\theta\;+\;\beta\,b\,\frac{D(v)}{D(v)^2+E(v)^2}.
\label{eq:thetav}
\end{equation}
Consequently the effect is point identified, $\theta(v)=\theta$, if and only if
$\beta b=0$ or $D(v)=0$, and the set of directions for which it is identified is
the subspace $\mathcal{V}_0=\{v:D(v)=0\}$ orthogonal to the missed spatial channel.
The residual-overlap quantity that Assumption~\ref{as:id-overlap} requires to be
positive is $\mathbb{E}[\tilde T^2]=D(v)^2+E(v)^2$, so a direction lying entirely
in the smooth spatial subspace violates positivity and leaves $\theta(v)$
undefined.
\end{theorem}

\begin{proof}[Proof]
Conditioning on $L$ removes the smooth channel, so the residualized treatment is
$\tilde T=D(v)\,\ell_d+E(v)\,\eta$ in the orthonormal channel basis, where $\ell_d$
is the missed spatial channel and $\eta$ the semantic one. The residualized
outcome is $\tilde Y=\theta\tilde T+\beta\,b\,\ell_d$, since $U$ residualized on
$L$ retains only its weight $b$ on $\ell_d$. Taking the Robinson ratio,
\[
\theta(v)=\frac{\mathbb{E}[\tilde Y\tilde T]}{\mathbb{E}[\tilde T^2]}
=\frac{\theta\,(D(v)^2+E(v)^2)+\beta b\,D(v)}{D(v)^2+E(v)^2}
=\theta+\beta b\,\frac{D(v)}{D(v)^2+E(v)^2},
\]
which is \eqref{eq:thetav}. The identification and positivity claims read off
directly. The algebra is verified symbolically over a free-symbol channel model,
and the estimand \eqref{eq:thetav} is recovered to within Monte Carlo error by a
cross-fitted gradient-boosting estimator across a sweep of directions, including
exact recovery of $\theta$ on $\mathcal{V}_0$.
\end{proof}

Three consequences deserve statement. The scalar identification result is the
special case of a single unit-variance semantic channel, $D(v)=d$ and $E(v)=1$,
where \eqref{eq:thetav} collapses to the bias $\beta bd/(d^2+1)$ of
Theorem~\ref{thm:id}; the vector result is therefore a strict generalization
rather than an analogy. The bias in \eqref{eq:thetav} is largest in magnitude
where $|D(v)|=|E(v)|$, a direction that mixes the missed spatial channel and the
semantic one in equal measure, and it decays as the direction tilts toward either
the identified subspace or the smooth channel that positivity forbids. And because
$\theta(v)$ varies continuously with $v$, a representation whose leading variance
direction happens to load on the missed spatial channel will report an effect that
is mostly laundered geography, while a direction chosen to maximize $E(v)$ relative
to $D(v)$ reports the semantic effect the analyst intends. Identification here is
not a property the embedding either has or lacks. It is a property of the direction
drawn through it, and the leading principal component, chosen for variance rather
than for identification, need not lie in the identified subspace.
Section~\ref{sec:identified-direction} carries this prescription out on the
twelve-market panel, re-estimating the effect on the identified direction and
reading the difference from the leading component as the geography it removes.

\paragraph{The estimand in estimable form.} Equation~\eqref{eq:thetav} is written in
the channel geometry that makes identification transparent, but the same estimand has
a form that depends only on quantities the data deliver. Writing $a=\mathbb{E}[\tilde
W\tilde Y]$ for the residualized embedding--price covariance vector and
$\Sigma=\mathbb{E}[\tilde W\tilde W']$ for the residualized embedding covariance, the
Robinson estimand on the direction $v$ is
\begin{equation}
\theta(v)=\frac{v'a}{v'\Sigma v},
\label{eq:thetav-estimable}
\end{equation}
of which \eqref{eq:thetav} is the two-channel special case. The identified subspace
$\mathcal{V}_0$ is the set of directions along which the numerator carries only the
semantic covariance, and positivity is $v'\Sigma v>0$. This is the object the panel
estimates, market by market, once $a$ and $\Sigma$ are replaced by their cross-fitted
sample analogues and $v$ by the identified or the leading direction.

\begin{corollary}[The variance-maximal direction can reverse the sign]
\label{cor:signflip}
Let the residualized embedding covariance place more variance on the missed spatial
channel than on the semantic one, $\sigma_s^2>\sigma_q^2$. Then the leading principal
component, the direction that maximizes $v'\Sigma v$, aligns with the spatial channel,
and $\theta(v_{\mathrm{PC}})\to b$, the price coefficient of the confounder itself
rather than the semantic effect $\theta$. When
$\operatorname{sign}(b)\neq\operatorname{sign}(\theta)$ the reported effect is
sign-reversed: the default summary of the text returns a coefficient of the wrong
sign.
\end{corollary}

The corollary is the theory behind New York. Its leading text direction is the most
geographic in the panel, so it loads on the borough and neighborhood channel whose
price gradient runs opposite the semantic effect, and the naive estimate on that
direction is a significant negative while the identified estimate is a significant
positive. \citet{guiveitch2023} observe a naive-versus-adjusted sign reversal
empirically in a different text setting, a politeness effect that moves from $-0.038$
to $+0.200$ under adjustment, but state no condition for it; Corollary~\ref{cor:signflip}
supplies the condition and identifies the variance-maximal direction, the one an
analyst reaches for by default, as the case at risk.

\paragraph{A spectral reading.} When the field is stationary with matrix spectral
density $\Sigma_W(\omega)$ and the confounder concentrates its mass on a
low-frequency band $\Omega_0$, the projected quantities become frequency integrals,
$D(v)^2+E(v)^2=\int v'\Sigma_W(\omega)v\,d\omega$ and the confounding numerator
$D(v)=\int_{\Omega_0} v'b(\omega)\,d\omega$ along the cross-spectrum $b(\omega)$
between the confounder and the field. Identification then asks that the direction
place energy outside the confounder's band, and the bias is bounded by a
Cauchy--Schwarz alignment between the direction and the residual confounding mass
that survives the control. This is the sense in which a text direction is
identified when it carries the field's high-frequency, non-spatial content and
confounded when it tracks the low-frequency geography the location control cannot
fully absorb.
